%% ============================================================================
%%  lafrox-portada.tex -- Caratula de oferta LafroX, variante "banda de cabecera"
%%
%%  Se carga automaticamente desde lafrox.cls. Define:
%%      \portadalafrox        caratula completa (folio 1)
%%      \contraportadalafrox  cierre opcional con el equipo y la declaracion
%%
%%  QUE CAMBIO EN LA v2.2 Y POR QUE
%%  -------------------------------
%%  Hasta la v2.1 la portada era una cuna diagonal a sangre que recorria la
%%  pagina completa y sobre la cual iba TODO el texto en negativo. Se veia
%%  bien, pero fallaba como caratula de licitacion por tres razones:
%%
%%   1. Jerarquia invertida. Los 38 pt se los llevaba "PROPUESTA TECNICA"
%%      --la etiqueta mas generica posible, todos los oferentes entregan
%%      una-- y el identificador de la licitacion iba en 9,5 pt gris. La
%%      caratula es un instrumento de IDENTIFICACION, no una tapa de marca:
%%      manda el numero de licitacion, el objeto, el sobre y el proponente.
%%
%%   2. No se puede firmar sobre tinta casi negra. El Art. 40.2 exige firma
%%      COMPLETA en la caratula, y el Sobre N. 1 se entrega fisico, foliado
%%      y firmado. Una firma de lapiz sobre #141414 es ilegible. Por eso el
%%      tercio inferior de la pagina queda ahora en blanco, sin excepcion:
%%      es zona de firma y de folio, no zona de arte.
%%
%%   3. El folio dependia de la geometria del arte. Iba en blanco sobre la
%%      diagonal exterior, que alcanzaba el margen derecho por 0,02 de ancho
%%      de papel. La falta de foliacion es causal de EXCLUSION (Art. 40.1):
%%      no puede depender de un parametro de diseno.
%%
%%  La banda diagonal se conserva, pero como cabecera: mismo lenguaje visual
%%  (diagonal, escalones de gris, monocroma estricta), una fraccion del area.
%%
%%  BLOQUE DEL ART. 49.1
%%  --------------------
%%  Las Bases fijan literalmente lo que debe leerse en la caratula:
%%      LICITACION N. TFEP-01/2026
%%      Proyecto de Plataforma Digital de Mision Critica -- Caso [N. y nombre]
%%      SOBRE N. x: [contenido]
%%      [Nombre del Proponente]
%%      [Nombre del Representante Legal]
%%      [Correo electronico y telefono del Representante]
%%  Los seis elementos estan abajo, en ese orden. Si se quita uno, la
%%  caratula deja de cumplir. Las claves que los alimentan (objeto, sobre,
%%  correo, telefono) se definen en lafrox.cls seccion 6.
%%
%%  Monocroma estricta. Ningun matiz: negro, grises neutros y blanco.
%% ============================================================================

% ---- Parametros de la banda de cabecera (fracciones del papel) -------------
%  La banda cruza el ancho completo y su borde inferior es diagonal: mas baja
%  a la derecha que a la izquierda. Para abrir o cerrar la diagonal basta
%  cambiar estos dos valores.
%
%  NO subir \LFXbandaIzqY por encima de ~0.19: el bloque de identificacion
%  empieza en 0.228 y se montaria sobre los escalones de gris.
\def\LFXbandaIzqY{0.150}    % altura de la banda en el borde izquierdo
\def\LFXbandaDerY{0.088}    % altura de la banda en el borde derecho

% ---- Utilidades de posicionamiento ----------------------------------------
%  \LFXen{fraccion x}{fraccion y}  -> coordenada respecto del vertice sup. izq.
%  Las llaves y el * son indispensables: sin ellas, una expresion como
%  0.17+0.013 se lee como "0.17 pt mas 0.013 anchos de papel".
\newcommand{\LFXen}[2]{%
  ([xshift={(#1)*\paperwidth},yshift={-(#2)*\paperheight}]current page.north west)}

% ---- Reticula horizontal de la portada ------------------------------------
%  \LFXmi y \LFXmd son las fracciones de papel que corresponden a los
%  margenes REALES del bloque de texto. Se calculan, no se escriben a mano:
%  asi la caratula cae exactamente sobre los mismos margenes que cualquier
%  pagina de contenido, y sigue cayendo ahi si manana se cambia \geometry.
%
%  La banda NO usa estas fracciones: es arte a sangre y va de borde a borde a
%  proposito.
\newcommand{\LFXcalcularmargenes}{%
  \pgfmathsetmacro{\LFXmi}{(1in+\oddsidemargin+\hoffset)/\paperwidth}%
  \pgfmathsetmacro{\LFXmd}{(1in+\oddsidemargin+\hoffset+\textwidth)/\paperwidth}%
}

% ---- Estado del documento, derivado de la opcion de clase ------------------
%  Va en el bloque de control documental y repetido en la banda. Un documento
%  de licitacion declara de si mismo en que estado esta; eso es justamente lo
%  que lo hace parecer --y ser-- un documento de licitacion.
\ifLFX@borrador
  \newcommand{\LFXestadodoc}{Borrador de trabajo}
\else
  \newcommand{\LFXestadodoc}{Versión final para entrega}
\fi

% ---- Marca en negativo: va dentro de la banda oscura -----------------------
%  ATENCION: si se define la clave \LFXlogo, el archivo debe ser una version
%  en NEGATIVO (trazo claro). Un logo normal, de trazo oscuro, desaparece
%  sobre la banda.
%  Sin \LFXlogo, la marca es el isotipo NARANJA del kit (logo/) + razon
%  social en blanco; si falta el archivo, vuelve al monograma provisional.
\newcommand{\LFXmarcaportada}{%
  \ifx\LFXlogo\@empty
    \begin{tikzpicture}[baseline=0pt]
      \node[anchor=west,inner sep=0pt] (mono) at (0,0)
        {\IfFileExists{\LFXisotipoportada}
           {\includegraphics[height=15mm]{\LFXisotipoportada}}
           {\lafroxmonogramaneg[11mm]}};
      \node[anchor=west,inner sep=0pt,align=left] at ([xshift=3.4mm]mono.east)
        {\sffamily\bfseries\fontsize{17pt}{17pt}\selectfont
         \color{white}\LFXempresa};
    \end{tikzpicture}%
  \else
    \includegraphics[height=13mm]{\LFXlogo}%
  \fi}

% ---- Campo rotulado: rotulo pequeno arriba, valor debajo -------------------
%  Es la unidad de composicion de los bloques de identificacion y de control.
%  \begin{tabular}[t] y no minipage: una minipage de ancho natural no existe,
%  y aqui los campos se reparten con \hfill segun su ancho real.
\newcommand{\LFXcampoportada}[2]{%
  \begin{tabular}[t]{@{}l@{}}
    {\color{lafroxGris}\LFXrotulo{#1}}\\[4pt]
    {\fontsize{10pt}{13.5pt}\selectfont\color{lafroxTinta}#2}
  \end{tabular}}

% ---- Filete separador de ancho de caja -------------------------------------
\newcommand{\LFXfileteportada}[1]{%
  \draw[lafroxLinea,line width=0.5pt] \LFXen{\LFXmi}{#1} -- \LFXen{\LFXmd}{#1};}

% ============================================================================
%  \portadalafrox
% ============================================================================
\newcommand{\portadalafrox}{%
  \begingroup
  \setcounter{page}{1}%
  \thispagestyle{lfxportada}%
  \LFXcalcularmargenes
  \null
  \begin{tikzpicture}[remember picture,overlay]

    % --- Banda de cabecera -------------------------------------------------
    %  El resto del papel queda en blanco a proposito. No agregar fondo
    %  \lafroxBg: el bloque de firma tiene que quedar sobre papel limpio.
    \fill[lafroxBanda]
      \LFXen{0}{0}              --
      \LFXen{1}{0}              --
      \LFXen{1}{\LFXbandaDerY}  --
      \LFXen{0}{\LFXbandaIzqY}  -- cycle;

    % --- Escalones de gris que acompanan el filo de la banda ---------------
    \fill[lafroxNaranja]
      \LFXen{0}{\LFXbandaIzqY}        --
      \LFXen{1}{\LFXbandaDerY}        --
      \LFXen{1}{\LFXbandaDerY+0.013}  --
      \LFXen{0}{\LFXbandaIzqY+0.013}  -- cycle;
    \fill[lafroxNaranjaC]
      \LFXen{0}{\LFXbandaIzqY+0.022}  --
      \LFXen{1}{\LFXbandaDerY+0.022}  --
      \LFXen{1}{\LFXbandaDerY+0.028}  --
      \LFXen{0}{\LFXbandaIzqY+0.028}  -- cycle;

    % --- Marca y estado, dentro de la banda --------------------------------
    \node[anchor=north west,inner sep=0pt] at \LFXen{\LFXmi}{0.050}
      {\LFXmarcaportada};
    \node[anchor=north east,inner sep=0pt] at \LFXen{\LFXmd}{0.056}
      {\color{lafroxSobre}\LFXrotulo{\LFXestadodoc}};

    % ======================================================================
    %  1-3. Licitacion, objeto y caso  (Art. 49.1, lineas 1 y 2)
    %  Aqui manda el identificador del proceso: es el dato con el que el
    %  evaluador ordena las ofertas sobre el escritorio.
    % ======================================================================
    \node[anchor=north west,inner sep=0pt] at \LFXen{\LFXmi}{0.228}
      {\color{lafroxGris}\LFXrotulo{Licitación pública}};
    \node[anchor=north west,inner sep=0pt] at \LFXen{\LFXmi}{0.248}
      {\sffamily\bfseries\fontsize{21pt}{25pt}\selectfont
       \color{lafroxFuerte}\LFXlicitacion};
    \node[anchor=north west,inner sep=0pt] at \LFXen{\LFXmi}{0.296}
      {\sffamily\fontsize{12.5pt}{16pt}\selectfont
       \color{lafroxTinta}\LFXobjeto};
    \node[anchor=north west,inner sep=0pt] at \LFXen{\LFXmi}{0.322}
      {\sffamily\fontsize{10.5pt}{14pt}\selectfont
       \color{lafroxGris}\LFXcaso};

    \LFXfileteportada{0.356}

    % ======================================================================
    %  4. Sobre y documento  (Art. 49.1, linea 3)
    %  El sello del sobre va en negativo: es el unico elemento de la mitad
    %  inferior con tinta plena, y marca de que pieza de la oferta se trata.
    % ======================================================================
    \node[anchor=north west,inner xsep=8pt,inner ysep=6pt,fill=lafroxBanda]
      at \LFXen{\LFXmi}{0.376}
      {\fontsize{8.5pt}{10pt}\selectfont
       \addfontfeature{LetterSpace=12}\intermedium
       \color{white}\LFXsobre};

    %  El titulo se toma de \LFXdocumento y NO va cableado: hasta la v2.1
    %  decia "PROPUESTA TECNICA" a mano, de modo que la portada mentia si la
    %  clase se usaba para la Oferta Economica o para un anexo.
    \node[anchor=north west,inner sep=0pt] at \LFXen{\LFXmi}{0.430}
      {\sffamily\bfseries\fontsize{31pt}{34pt}\selectfont
       \color{lafroxFuerte}\LFXdocumento};
    \node[anchor=north west,inner sep=0pt] at \LFXen{\LFXmi}{0.486}
      {\sffamily\fontsize{16pt}{20pt}\selectfont
       \color{lafroxTinta}\LFXsubtitulo};
    \node[anchor=north west,inner sep=0pt] at \LFXen{\LFXmi}{0.516}
      {\sffamily\fontsize{9.5pt}{13pt}\selectfont\color{lafroxGris}
       \LFXalcance\ \textperiodcentered\ \LFXformulario};

    \LFXfileteportada{0.552}

    % ======================================================================
    %  5-6. Mandante, proponente y representante  (Art. 49.1, lineas 4 a 6)
    %  El mandante se rotula como tal y no como parte de una frase: en una
    %  caratula real va con formalidad al menos igual que la del oferente.
    % ======================================================================
    \node[anchor=north west,inner sep=0pt] at \LFXen{\LFXmi}{0.570}
      {\begin{minipage}[t]{0.44\textwidth}
         \sffamily\raggedright
         {\color{lafroxGris}\LFXrotulo{Mandante}}\\[4pt]
         {\fontsize{11pt}{14pt}\selectfont\intersemibold
          \color{lafroxTinta}\LFXcliente}\\[13pt]
         {\color{lafroxGris}\LFXrotulo{Proponente}}\\[4pt]
         {\fontsize{11pt}{14pt}\selectfont\intersemibold
          \color{lafroxTinta}\LFXempresa}\\[1pt]
         {\fontsize{9.5pt}{13pt}\selectfont\color{lafroxTinta}
          RUT \LFXrut\\ \LFXdomicilio}
       \end{minipage}};

    \node[anchor=north west,inner sep=0pt] at \LFXen{\LFXmi+0.390}{0.570}
      {\begin{minipage}[t]{0.44\textwidth}
         \sffamily\raggedright
         {\color{lafroxGris}\LFXrotulo{Representante legal}}\\[4pt]
         {\fontsize{11pt}{14pt}\selectfont\intersemibold
          \color{lafroxTinta}\LFXrepresentante}\\[1pt]
         {\fontsize{9.5pt}{13pt}\selectfont\color{lafroxTinta}
          \LFXcargo\\ \LFXcorreo\\ \LFXtelefono}
       \end{minipage}};

    \LFXfileteportada{0.714}

    % ======================================================================
    %  Control documental
    %  El total de folios sale de \pageref{LastPage} (paquete lastpage):
    %  declararlo en la caratula permite detectar un PDF truncado antes de
    %  que lo detecte el evaluador, y la foliacion es excluyente (Art. 40.1).
    % ======================================================================
    \node[anchor=north west,inner sep=0pt] at \LFXen{\LFXmi}{0.732}
      {\begin{minipage}{\textwidth}
         \sffamily
         \LFXcampoportada{Fecha de emisión}{\LFXfecha}\hfill
       \end{minipage}};

    \LFXfileteportada{0.802}

    % ======================================================================
    %  Firma completa del representante  (Art. 40.2)
    %  Sobre papel blanco y con los datos BAJO la linea, para que la firma
    %  sea atribuible. No mover este bloque sobre la banda ni agregarle
    %  fondo: se firma con lapiz.
    % ======================================================================
    \node[anchor=north west,inner sep=0pt] at \LFXen{\LFXmi}{0.820}
      {\begin{minipage}[t]{0.52\textwidth}
         \sffamily\raggedright
         {\color{lafroxGris}\LFXrotulo{Firma del representante legal}}\\[34pt]
         {\color{lafroxTinta}\rule{6.4cm}{0.5pt}}\\[4pt]
         {\fontsize{9.5pt}{13pt}\selectfont\color{lafroxTinta}
          \intersemibold\LFXrepresentante\\
          \normalfont\sffamily\LFXcargo\ \textperiodcentered\ \LFXempresa}
       \end{minipage}};

    % --- Nomenclatura del archivo, a la altura del folio -------------------
    %  Art. 50.3: un archivo mal nominado se considera NO PRESENTADO. Va al
    %  pie izquierdo, emparejado con el folio del pie derecho, que es donde
    %  un codigo de archivo se busca. Antes tenia bloque propio y empujaba la
    %  firma sobre el pie.
    \node[anchor=base west,inner sep=0pt] at \LFXen{\LFXmi}{0.9565}
      {\ttfamily\fontsize{8pt}{10pt}\selectfont
       \color{lafroxGris}\LFXnomenclatura};

  \end{tikzpicture}
  \clearpage
  \endgroup}

% ============================================================================
%  \contraportadalafrox -- cierre institucional (opcional)
%  Uso: \contraportadalafrox{tabla de equipo o texto libre}
%
%  Mantiene la banda superior de la caratula, con la misma diagonal, para que
%  las dos tapas se lean como un par.
% ============================================================================
\newcommand{\contraportadalafrox}[1]{%
  \clearpage
  \thispagestyle{lfxcierre}%
  \LFXcalcularmargenes
  \null
  \begin{tikzpicture}[remember picture,overlay]
    \fill[lafroxBanda]
      \LFXen{0}{0}                   --
      \LFXen{1}{0}                   --
      \LFXen{1}{0.230}               --
      \LFXen{0}{0.296}               -- cycle;
    \fill[lafroxNaranja]
      \LFXen{0}{0.296}               --
      \LFXen{1}{0.230}               --
      \LFXen{1}{0.243}               --
      \LFXen{0}{0.309}               -- cycle;
    \node[anchor=north west,inner sep=0pt,align=left] at \LFXen{\LFXmi}{0.110}
      {\sffamily\bfseries\fontsize{26pt}{29pt}\selectfont\color{white}
       \LFXempresa};
    %  Dos nodos y no uno con \\: dentro de un nodo TikZ el salto de linea
    %  reinicia los ajustes de fuente Y DE COLOR, de modo que la segunda
    %  linea volvia a \lafroxTinta y quedaba negro sobre negro, ilegible.
    %  Es el mismo motivo por el que el titular de la portada usa dos nodos.
    \node[anchor=north west,inner sep=0pt] at \LFXen{\LFXmi}{0.190}
      {\sffamily\fontsize{10.5pt}{14pt}\selectfont\color{lafroxSobre}
       \LFXdocumento\ \textperiodcentered\ \LFXsubtitulo};
    \node[anchor=north west,inner sep=0pt] at \LFXen{\LFXmi}{0.212}
      {\sffamily\fontsize{10.5pt}{14pt}\selectfont\color{lafroxSobre}
       \LFXlicitacion\ \textperiodcentered\ \LFXcaso};
  \end{tikzpicture}
  \vspace*{0.26\paperheight}
  #1
  \clearpage}
